\section{Discussion}

The guarantees in \cref{obj:thm:adaptive-minimax} describe response recovery with known smoothness and adaptation to overlap. The known smoothness parameter \(\beta\) determines the polynomial degree and the approximation order used in constructing the estimator. Treatment counts then select a feasible mesh through \cref{obj:def:mesh-selector}. Thus response regularity supplies the approximation benchmark, while the observed treatment design determines the spatial resolution.

The compact range for the overlap parameter \(\gamma\), contained in \((1,\infty)\), specifies the domain over which the risk constants are uniform. With the remaining model constants fixed, a single estimator attains the corresponding oracle scale throughout that range. The pointwise and expected spatial-supremum guarantees therefore share an overlap-dependent scale, including the spatial-supremum guarantee without an additional logarithmic factor.

For comparison, the classical regular-design spatial-supremum benchmark includes a logarithmic factor, with scale \((\log n/n)^{\beta/(2\beta+d)}\) under its regular coverage conditions \citep{Stone1982}. The present theorem instead takes a worst case over the larger global-tail class at each fixed finite \(\gamma\), giving the pure power \(n^{-\beta/(2\beta+d\gamma/(\gamma-1))}\). The theorem supplies one estimator and common risk constants on every prescribed compact interval of finite overlap exponents.

Equal-cell weighting connects the deterministic geometry in \cref{obj:def:template} to the fitted polynomials in \cref{obj:def:estimator}. Each microcell contributes the same total weight, so the conditioning argument is governed by the template geometry, while differences in treatment counts govern stochastic precision; \cref{obj:lem:template-conditioning,obj:lem:ordered-mass} give the two bounds.

Conditional sub-Gaussian residuals supply exponential control of fluctuations after conditioning on the sampled covariates and treatment indicators. That control supports the upper-bound argument and its Gaussian-envelope transfer in \cref{obj:prop:dorn-a3-free-corollary}. Together, the response smoothness, template geometry, and treatment-count information determine approximation, conditioning, and stochastic error, respectively.

\subsection{Limitations and future work}

Joint adaptation would replace the known smoothness benchmark with a comparison across response regularities as well as treatment-count-feasible resolutions. The following proposal describes an architecture for that extension.
\begin{remarkv}{def:joint-adaptation-handle}[Future work: joint adaptation in $\beta$ and $\gamma$]
A proposed direction for future work is to construct a single response estimator with an adaptive-risk envelope over compact rectangles of the smoothness parameter \(\beta\) and the overlap parameter \(\gamma\). The proposed architecture uses sample splitting to compare equal-cell fits across polynomial degrees and treatment-count-feasible meshes, with the count selector controlling adaptation to overlap and held-out Lepski inequalities controlling adaptation to smoothness. The sample split, degree-and-mesh comparison rule, Lepski threshold, estimator, and envelope functional remain to be specified; estimator existence and the corresponding risk bounds remain open.
\end{remarkv}

The proposed split would allocate separate observations to fitting and comparison, while the degree-and-mesh rule would compare candidates along both parameter axes. Lepski comparisons provide a natural perspective on smoothness selection \citep{lepski1997}; smoothness adaptation under degenerate random design, studied by \citet{Gaiffas2005AdaptiveDegenerateDesign}, supplies a related setting in which to interpret the adaptation question. These comparisons motivate the proposal rather than establish its risk envelope.

The absence of an additional spatial-supremum logarithm in the known-smoothness guarantee makes the cost of the proposed two-axis comparison particularly relevant. The question concerns both attainability of the oracle scale and any unavoidable adaptation penalty.
\begin{remarkv}{oeq:joint-adaptation}[Joint smoothness and overlap adaptation]
For a compact rectangle of smoothness and overlap parameters \((\beta,\gamma)\) contained in \((1,\infty)^2\), can a procedure adapting along both parameter axes produce a single estimator whose pointwise and expected spatial-supremum risks attain the oracle scale for every pair in the rectangle? What logarithmic penalty, if any, is minimax unavoidable?
\end{remarkv}

Resolving \cref{obj:oeq:joint-adaptation} requires specifying the procedure and establishing its risk bounds, together with an appropriate converse for any claimed unavoidable logarithmic price. The existing common-class bounded-response converse and Gaussian-envelope upper transfer address distinct comparison domains: the former supplies the minimax lower bound in \cref{obj:thm:adaptive-minimax}, while the latter supplies the upper guarantee in \cref{obj:prop:dorn-a3-free-corollary}. The transferred upper bound does not furnish a matching converse for every Gaussian source subfamily. As \cref{obj:lem:arbitrary-family-lower-obstruction,obj:lem:dorn-rate-scope} clarify, a lower-bound analysis must account for family richness. These distinctions also govern how a future joint-adaptation result could characterize an unavoidable logarithmic penalty.
